% !TeX root = ../教材题库_全部条目_转写工作稿.tex
% Source: XJ-B1-C05-TH-P184-Y1253
% 第5章 / 5.3 三角函数的图象与性质 / 教材 p.184、185 / PDF 第191、192页
% OCR未校对：公式与文字均可能有误；未完成数学TeX转写；禁止直接发给学生
% 原题及相邻原解截图（核对用）：../题目原图/XJ-B1-C05-TH-P184-Y1253-01.webp, ../题目原图/XJ-B1-C05-TH-P184-Y1253-02.webp
% 本文件只在转写完成后作为题干使用；OCR 条目尚不是可靠题干。
\begin{bankitem}{p.184 思考（含紧邻说明）}{第5章 / 5.3 三角函数的图象与性质 / 教材 p.184、185 / PDF 第191、192页}{XJ-B1-C05-TH-P184-Y1253}{OCR未校对：公式与文字均可能有误；未完成数学TeX转写；禁止直接发给学生}
\begin{ocrdraft}
[教材p.184 / PDF第191页]
\par
？思考类比对正弦函数、余弦函数性质的研究，你能发现正切函数3，一tan工
\par
具有哪些性质？
\par
1.周期性
\par
由上面的0）式可知，正切函数3，一tan是周期函数，（差ez，差乒0）是它的周
\par
期，而丌是正切函数的最小正周期．
\par
2\ensuremath{\cdot}值域
\par
由图5．3一9可知，随着ZAOIP的弧度数x从0开始增大，正切线AT的长度
\par
也在增大，并且当，趋向于《时，，AT的长度可以大于指定的任意正数，即趋向于无
\par
[教材p.185 / PDF第192页]
\par
穷大，则tan工能取到@，十亠）内的任意实数．类似地，tan工也能取到（一。。，0]
\par
内的任意实数．因此，正切函数的值域是实数集R，正切函数没有最大值和最小值．
\par
3\ensuremath{\cdot}奇偶性
\par
观察正切曲线，我们可以看到正切曲线关于原点0对称，正切函数y=tan应
\par
该是奇函数．由诱导公式tan((x)=—tanx知，正切函数3，一tan是奇函数
\par
4\ensuremath{\cdot}单调性
\par
从图5，3一10可以看出：正切函数3一t“，在每个开区间（
\par
一十差丌，
\par
2
\par
（差ez）上单调递增．
\end{ocrdraft}
\end{bankitem}
